\section{Discussion and lessons learned}
\label{sec:discussion}

The evaluation separates document discovery from answer construction. On the
tested sets, retrieval-only recall does not by itself determine whether an
agent returns an addressable source, an explicit partial answer or a refusal.
The useful design lesson is to expose the retrieval trace and the evidence
returned to the answering step, while distinguishing a completed search with
no support from a search that did not complete.

The first implication concerns verifier placement. On identical logged pools,
re-ranking by the judge lowers Hit@5 from 0.94 under fusion-only to 0.44, while
a derived replay that keeps the fusion order among scored survivors returns
the fusion-only metrics. Because all 191 recorded scores exceed 0.5 and Q14
has one fused candidate without a recorded score, this does not demonstrate
threshold-filter benefit. It shows that the tested judge scores were a poor
sorting signal on this confuser-heavy corpus. The second verifier replay and
the second corpus do not establish universal calibration; the strongest
supported claim is that ranking by this checkpoint should not replace fusion
order under the tested conditions.

The second implication concerns abstention. Similarity can retrieve related
material without answering a query. In this system, the answer workflow
instructs the agent to report the searched scope and say not-found when its
completed evidence review finds no support. The adversarial battery observed
no fabricated citations in eighteen final arms, but one partial answer fell
outside the binary grading scheme and each arm has a single final observation.
Timeouts, unavailable tools and incomplete searches must be reported as
execution failures, not conflated with semantic absence.

The third implication concerns transport. A reproduced oversized-catalog
response hid entries from the agent and led to a false refusal. Server-side
budgets, grouped views and pagination repaired the selected tested paths on
rerun. This supports the narrower point that transport limits can affect
retrieval outcomes; it is not a guarantee for every host or larger catalog.

The five research questions close at different evidence levels. RQ1 is partly
answered: the system combines server-side parse/index checks and an invoked
content gate with agent procedures for placement and answer reporting. RQ2 is
answered as an operational observation, bounded by 21 papers that never became
retrievable and the absence of independent placement ground truth. RQ3 reports
tool-layer recall separately from end-to-end source-marker checks; it does not
claim similar recall in every interface or isolate the gate's causal effect.
RQ4 identifies a same-pool ranking failure for one checkpoint; it does not
validate threshold discrimination. RQ5 documents observed abstentions,
turn-dominated latency and a transport failure/recovery, with no broad
reliability or user-utility estimate.

More broadly, parsers, registries, indices and answer workflows benefit from
post-conditions that can be checked independently of success messages.
Catalogues and graphs support navigation but are not themselves evidence, in
keeping with provenance and auditability work in data systems
\cite{missier-prov-2013,cheney-provenance-2009}. The results remain bounded by
the tested corpora, developer-authored questions, one workstation and few
end-to-end observations. Larger independent question sets, repeated agent
trials, claim-level support assessments, operator studies of not-found
reports and crash-recovery tests are needed before stronger deployment claims.
